\documentclass[../../../include/open-logic-section]{subfiles}

\begin{document}

\olfileid{sth}{choice}{vitali}
\olsection{परिशिष्ट: वितालीचा विरोधाभास}

बानाख--टार्स्कीची रचना स्वीकारार्ह आहे की नाही, हे नीट समजून घेण्यासाठी तिचा \emph{पुरावा} तपासला पाहिजे. दुर्दैवाने त्यासाठी येथे मांडता येईल त्याहून बरेच अधिक बीजगणित लागेल. तरी पुढील निष्कर्षावर लक्ष केंद्रित करून बानाख--टार्स्कीच्या पुराव्यावर प्रकाश टाकतील अशा काही थोडक्यात नोंदी देता येतील.\footnote{अधिक पूर्ण विवेचनासाठी \cite{Weston2003} किंवा \cite{Wagon2016} पाहा.}

\begin{thm}[वितालीचा विरोधाभास ($\ZFC$ मध्ये)]\ollabel{vitaliparadox}
धन त्रिज्येच्या कोणत्याही वर्तुळाचे गणनीय संख्येने तुकडे करता येतात; घूर्णन आणि स्थानांतरणाने त्यांची पुनर्जुळणी करून त्या वर्तुळाच्या दोन प्रती तयार करता येतात.
\end{thm}

वितालीचा विरोधाभास बानाख--टार्स्की विरोधाभासापेक्षा सिद्ध करणे खूप सोपे आहे. वितालीने \citeyear{Vitali1905} मध्ये केलेल्या अमापनीय संचाच्या रचनेवरून तो मिळत असल्यामुळे त्याला “वितालीचा विरोधाभास” म्हटले आहे. पण विताली आणि बानाख--टार्स्की यांच्या पुराव्यांचे संचसैद्धान्तिक पैलू अगदी सारखे आहेत. मुख्य फरक एवढाच की बानाख--टार्स्कीमध्ये \emph{सांत} विघटन आहे, तर वितालीमध्ये \emph{गणनीय अनंत} विघटन आहे. \citet{Weston2003} यांच्या शब्दांत, वितालीचा विरोधाभास “बानाख--टार्स्की विरोधाभासाइतका नक्कीच थक्क करणारा नाही; पण अगदी ‘साध्या’ परिस्थितीतही भूमितीय विरोधाभास येऊ शकतात, हे तो दाखवतो.”

वितालीचा विरोधाभास वर्तुळ या द्विमितीय आकृतीविषयी आहे. म्हणून $\Real^2$ या प्रतलावर काम करू. मूलबिंदूभोवती पूर्ण फेरीच्या \emph{परिमेय} भागांनुसार, $[0,2\pi)$ या मर्यादेत रेडियनमूल्य असलेली दक्षिणावर्त घूर्णने $\rotationsgroup$ मध्ये घ्या. येथे पूर्ण फेरीचा परिमेय भाग अभिप्रेत आहे; परिमेय संख्येइतके रेडियन नव्हेत. $\rotationsgroup$ विषयी पुढील बीजगणितीय तथ्ये पाहा (विधान समजले नाही, तर पुराव्यात त्याचा अर्थ स्पष्ट होईल):

\begin{lem}\ollabel{rotationsgroupabelian}
फलनांच्या संयोजनाखाली $\rotationsgroup$ हा अबेली, म्हणजे क्रमनिरपेक्ष, {गट} होतो.
\end{lem}

\begin{proof}
$0$ रेडियनचे घूर्णन $0_{\rotationsgroup}$ असे लिहिल्यास हा $\rotationsgroup$ चा अविकारी घटक आहे. कारण प्रत्येक $\rho \in \rotationsgroup$ साठी $\comp{0_{\rotationsgroup}}{\rho} = \comp{\rho}{0_{\rotationsgroup}} =
\rho$. पूर्ण फेऱ्या वगळून कोन घेतल्यास घूर्णनांचे संयोजनही याच गटात राहते.

प्रत्येक घटकाला व्यस्त असतो. शून्य घूर्णनाचा व्यस्त तोच असतो. इतर प्रकरणांत, $\rho \in \rotationsgroup$ हे $r$ रेडियनचे घूर्णन असेल, तर $\rho^{-1} \in \rotationsgroup$ हे $2\pi - r$ रेडियनचे घूर्णन असते. त्यामुळे $\rho \circ \rho^{-1} = 0_\rotationsgroup$.

संयोजनात सहचारिता आहे: कोणत्याही $\rho, \sigma, \tau \in
\rotationsgroup$ साठी $\comp{\rho}{(\comp{\sigma}{\tau})} =
\comp{(\comp{\rho}{\sigma})}{\tau}$.

संयोजन क्रमनिरपेक्ष आहे: कोणत्याही $\rho, \sigma \in \rotationsgroup$ साठी $\comp{\rho}{\sigma} =
\comp{\sigma}{\rho}$.
\end{proof}

खरे तर $\rotationsgroup$ या गटाचे दोन भाग करून प्रत्येक भागावरून पूर्ण गट पुन्हा मिळवता येतो:

\begin{lem}\ollabel{disjointgroup}
$\rotationsgroup$ चे $\rotationsgroup_{1}$ आणि $\rotationsgroup_{2}$ अशा दोन वियुक्त संचांत विभाजन करता येते; दोन्ही $\rotationsgroup$ चे जनक संच आहेत, आणि पुढील पुराव्यात त्यांच्या दुप्पट घूर्णनांची प्रतिमा पूर्ण गट देते, हे अधिक नेमके दाखवले आहे.
\end{lem}

\begin{proof}
पूर्ण फेरीच्या परिमेय भागांशी जुळणारी आणि रेडियनमूल्य $[0, \pi)$ मध्ये असलेली घूर्णने $\rotationsgroup_{1}$ मध्ये घ्या; आणि $\rotationsgroup_{2} = \rotationsgroup
\setminus \rotationsgroup_1$ घ्या. प्राथमिक बीजगणिताने
$\Setabs{\comp{\rho}{\rho}}{\rho \in \rotationsgroup_1} =
\rotationsgroup$. $\rotationsgroup_2$ साठीही तसाच निष्कर्ष मिळतो. प्रत्येक अर्ध्यात कोन दुप्पट करण्याचे प्रतिचित्रण पूर्ण गटावर एकास-एक आणि आच्छादक असते.
\end{proof}

गटांविषयीचे हे तथ्य \olref{vitaliparadox} सिद्ध करण्यासाठी वापरू. $\onesphere$ हे एकक वर्तुळ, म्हणजे प्रतलाच्या मूलबिंदूपासून नेमके $1$ एकक अंतरावरील बिंदूंचा संच, म्हणजे $\Setabs{\tuple{r,s} \in \Real^2}{\sqrt{r^2+s^2}=1}$ असू द्या. $\onesphere$ वरील पुढील संबंध पाहून $\onesphere$ चे भाग करू:
\[
	r \sim s \emph{ असल्यास आणि तरच }(\exists \rho \in \rotationsgroup)\rho(r) = s.
\]
म्हणजे मूलबिंदूभोवती पूर्ण फेरीच्या परिमेय भागाचे घूर्णन करून एका बिंदूपासून दुसऱ्याकडे जाता येत असल्यास आणि तरच $\onesphere$ चे ते बिंदू या संबंधाने जोडलेले आहेत. अपेक्षेप्रमाणे:

\begin{lem}
$\sim$ हा तुल्यता संबंध आहे.
\end{lem}

\begin{proof}
\olref{rotationsgroupabelian} वापरून हे तात्काळ मिळते.
\end{proof}

आता निवड वापरून $\sim$ नुसार $\onesphere$ च्या प्रत्येक तुल्यतावर्गातून नेमका एक घटक असलेला संच $C$ मिळवू. म्हणजे तुल्यतावर्गांच्या पुढील संचावर निवडफलन $f$ घेऊ:\footnote{$\rotationsgroup$ हे !!{enumerable} असल्यामुळे $E$ चा प्रत्येक !!{element} !!{enumerable} आहे. $\onesphere$ हे !!{nonenumerable} असल्यामुळे \olref{vitalicover} आणि \olref[card-arithmetic][simp]{kappaunionkappasize} नुसार $E$ हे !!{nonenumerable} आहे. म्हणून येथे \emph{अ}गणनीय निवड वापरली आहे.}
\[
	E = \Setabs{\equivrep{r}{\sim}}{r \in \onesphere},
\]
आणि $C = \ran{f}$ घ्या. प्रत्येक घूर्णन $\rho \in \rotationsgroup$ साठी, $C$ च्या प्रत्येक बिंदूवर $\rho$ लागू करून मिळणाऱ्या बिंदूंचा संच $\funimage{\rho}{C}$ आहे. पुढील दोन निष्कर्ष दाखवतात की हे संच वर्तुळ पूर्णपणे व्यापतात आणि एकमेकांवर आच्छादित होत नाहीत:

\begin{lem}\ollabel{vitalicover}
$\onesphere = \bigcup_{\rho \in \rotationsgroup} \funimage{\rho}{C}$.
\end{lem}

\begin{proof}
$s \in \onesphere$ निश्चित करा. $r \in \equivrep{s}{\sim}$ असलेला $r \in C$ आहे; म्हणजे $r \sim s$, म्हणजे एखाद्या $\rho \in \rotationsgroup$ साठी $\rho(r) = s$.
\end{proof}

\begin{lem}\ollabel{vitalinooverlap}
$\rho_1 \neq \rho_2$ असल्यास $\funimage{\rho_1}{C} \cap \funimage{\rho_2}{C} = \emptyset$.
\end{lem}

\begin{proof}
$s \in \funimage{\rho_{1}}{C} \cap \funimage{\rho_{2}}{C}$ असे समजा. मग एखाद्या $r_{1}, r_{2} \in C$ साठी $s = \rho_{1}(r_{1}) = \rho_{2}(r_{2})$.
त्यामुळे $\rho^{-1}_2(\rho_1(r_1)) = r_2$, आणि $\comp{\rho_1}{\rho^{-1}_2} \in \rotationsgroup$; म्हणून $r_{1} \sim
r_{2}$. $\sim$ नुसार प्रत्येक तुल्यतावर्गातून $C$ नेमका एक घटक निवडतो, म्हणून $r_1 = r_2$. त्यामुळे $s = \rho_1(r_1) = \rho_2(r_1)$ आणि म्हणून $\rho_1 = \rho_2$. येथे एकक वर्तुळाचा बिंदू मूलबिंदू नसतो; अशा बिंदूची घूर्णनाखालील प्रतिमा त्या घूर्णनाला नेमके ठरवते, हे वापरले आहे.
\end{proof}

आता आधीची बीजगणितीय तथ्ये वर्तुळाला लागू करू:

\begin{lem}\ollabel{pseudobanachtarski}
$\onesphere$ चे $D_{1}$ आणि $D_{2}$ अशा दोन वियुक्त संचांत विभाजन करता येते, ज्यामध्ये $D_{1}$ चे गणनीय संख्येने भाग करून ते घूर्णित केल्यास $\onesphere$ ची प्रत तयार होते (आणि $D_{2}$ साठीही तसेच).
\end{lem}

\begin{proof}
\olref{disjointgroup} मधील $\rotationsgroup_{1}$ आणि $\rotationsgroup_{2}$ वापरून पुढील व्याख्या करा:
\begin{align*}
	D_{1} &= \bigcup_{\rho \in \rotationsgroup_1} \funimage{\rho}{C} &
	D_{2} &= \bigcup_{\rho \in \rotationsgroup_2} \funimage{\rho}{C}
\end{align*}
\olref{vitalicover} नुसार हे $\onesphere$ चे विभाजन आहे, आणि \olref{vitalinooverlap} नुसार $D_1$ आणि $D_2$ वियुक्त आहेत. रचनेनुसार $D_1$ चे गणनीय संख्येने भाग करता येतात: प्रत्येक $\rho \in \rotationsgroup_1$ साठी $\funimage{\rho}{C}$ हा एक भाग. \olref{disjointgroup} आणि \olref{vitalicover} नुसार $\onesphere = \bigcup_{\rho \in
\rotationsgroup}\funimage{\rho}{C} = \bigcup_{\rho \in
\rotationsgroup_1}\funimage{(\comp{\rho}{\rho})}{C}$ असल्यामुळे हे भाग घूर्णित करून $\onesphere$ ची प्रत तयार करता येते. कोन दुप्पट करण्याचे प्रतिचित्रण प्रत्येक अर्ध्यावर एकास-एक असल्याने परिणामी भागही वियुक्त असतात. $D_2$ साठीही तोच युक्तिवाद लागू होतो. \end{proof}\noindent यावरून वितालीचा विरोधाभास लगेच मिळतो. वर्तुळाचे गणनीय संख्येने तुकडे (विविध $\funimage{\rho}{C}$) करून दृढ गतीने ते हलवल्यास $\onesphere$ पासून $\onesphere$ च्या \emph{दोन} प्रती तयार करता येतात. फक्त $D_1$ चा प्रत्येक भाग घूर्णित करा; आणि $D_2$ च्या प्रत्येक भागाचे आधी दुसरीकडे स्थानांतरण करून नव्या केंद्राभोवती घूर्णन करा.

पुराव्याची पद्धत पुन्हा पाहू. प्रतलावरील घूर्णन गटाच्या काही बीजगणितीय तथ्यांपासून सुरुवात केली. हा गट वापरून $\onesphere$ चे तुल्यतावर्गांत विभाजन केले. त्यानंतर प्रत्येक वर्गातून घटक निवडण्यासाठी निवड वापरून “विरोधाभास” मिळवला.

बानाख--टार्स्की सिद्ध करण्यासाठी नेमकी हीच पद्धत वापरतो. मुख्य फरक असा की त्यासाठीची बीजगणितीय तथ्ये वितालीच्या पुराव्यातील तथ्यांपेक्षा बरीच गुंतागुंतीची आहेत. पण त्या बीजगणितीय तथ्यांचा निवडीशी संबंध नाही. त्यांचा थोडक्यात सारांश देऊ.

बानाख--टार्स्की सिद्ध करण्यासाठी आधी \olref{disjointgroup} सारखा निष्कर्ष मिळवतो: दोन जनकांवरील \emph{मुक्त गटाचे} चार भाग करता येतात, आणि सहज अर्थाने ते “हलवून” पूर्ण गटाच्या दोन प्रती मिळवता येतात.\footnote{\emph{चार} भाग वापरता येतात, हा निष्कर्ष \cite{Robinson1947} यांचा आहे. अलीकडील पुराव्यासाठी \citet[Theorem 5.2]{Wagon2016} पाहा. गटाचे भाग “हलवणे” असे म्हणताना \citet[p.~3]{Weston2003} यांच्या वर्णनाचा आधार घेतला आहे.} त्यानंतर $\Real^3$ च्या मूलबिंदूभोवतीची दोन विशिष्ट घूर्णने वापरून घूर्णनांचा मुक्त गट $F$ तयार करता येतो, हे दाखवतो.\footnote{पाहा \citet[Theorem 2.1]{Wagon2016}.} (अजून निवड वापरलेली नाही.) आता $F$ मधील घूर्णनाने एका बिंदूपासून दुसरा मिळत असल्यास आणि तरच गोलकाच्या पृष्ठभागावरील ते बिंदू “सदृश” मानतो. मग “सदृश” बिंदूंच्या प्रत्येक तुल्यतावर्गातून नेमका एक बिंदू निवडण्यासाठी \emph{निवड वापरतो}. \olref{pseudobanachtarski} प्रमाणे $F$ चे विभाजन पृष्ठभागाला लागू करण्याची कल्पना अशी की चार भाग “हलवून” पृष्ठभागाच्या दोन प्रती मिळवाव्यात. मात्र या त्रिमितीय रचनेत अक्षांवरील काही बिंदू घूर्णनांनी स्थिर राहतात; अशा बिंदूंचा स्वतंत्र विचार आणि आवश्यक अतिरिक्त युक्तिवाद या संक्षिप्त रूपरेषेत दिलेला नाही. म्हणून ही पूर्ण पुराव्याची जागा घेऊ शकत नाही. संबंधित निष्कर्ष \citep{Hausdorff1914} असा मांडला जातो:

\begin{thm}[हाउसडॉर्फचा विरोधाभास ($\ZFC$ मध्ये)]
धन त्रिज्येच्या कोणत्याही गोलकाच्या पृष्ठभागाचे सांत संख्येने तुकडे करता येतात; घूर्णन आणि स्थानांतरणाने त्यांची पुनर्जुळणी करून त्या पृष्ठभागाच्या दोन वियुक्त प्रती तयार करता येतात.
\end{thm}

पूर्ण बानाख--टार्स्की प्रमेय मिळवण्यासाठी (जे फक्त गोलकाच्या पृष्ठभागाविषयी नाही, तर अंतर्भागाविषयीही आहे) आणखी काही बीजगणितीय युक्त्या लागतात. तरी लेखकाच्या मते मुख्य बीजगणितीय काम आधीच्या टप्प्यावर आहे. म्हणून वेस्टन लिहितात:
\begin{quote}
  [\ldots] मुक्त गटांविषयीचा निष्कर्ष बानाख--टार्स्की विरोधाभासाच्या पुराव्यातील \emph{मुख्य टप्पा} आहे. या दृष्टीने बानाख--टार्स्की विरोधाभास हा $\Real^3$ विषयीच्या विधानापेक्षा [$\Real^3$ मधील स्थानांतरणे आणि घूर्णने यांच्या] गटाच्या गुंतागुंतीविषयीचे विधान आहे. \cite[p.~16]{Weston2003}
\end{quote}
म्हणजे बानाख--टार्स्कीप्रमाणे \emph{सांत} विघटन मिळते की वितालीप्रमाणे \emph{गणनीय अनंत} विघटन मिळते, हे दोन किंवा तीन मितींत काम करताना गटांविषयी लागू होणाऱ्या काही तथ्यांवर ठरते.

हे शेवटचे निरीक्षण \olref[banach]{sec} च्या अखेरचा विनोद थोडा बिघडवते. लेखन द्विमितीय असल्यामुळे “बानाख--टार्स्की बानाख--टार्स्की” अशी पुनर्जुळणी करायची असेल, तर “बानाख--टार्स्की”चे गणनीय \emph{अनंत} तुकडे करावे लागतील. विनोद दुरुस्त करण्यासाठी त्रिमितीय लेखन करावे. हा सराव वाचकासाठी सोडला आहे.

शेवटची एक नोंद. \olref[banach]{sec} मध्ये गोलकाचे मिळणारे “तुकडे” सर्व \emph{मापनीय} असू शकत नाहीत; त्यात चित्रात दाखवता न येणारे “अनंत विखुरणे” असते, असे म्हटले होते. \olref{pseudobanachtarski} मिळवण्यासाठी केलेल्या निवडीच्या वापराविषयीही हेच लागू होते. हे स्पष्ट करणे उपयुक्त आहे.

पुन्हा थोडी पार्श्वभूमी सांगावी लागेल (ही \emph{केवळ} रूपरेषा आहे; \emph{मापाविषयी} पाठ्यपुस्तकातील विवेचन पाहावे). येथे वापरल्या जाणाऱ्या सांत मापात, संच $X$ वरील एखाद्या “$\sigma$-बीजगणितातील” प्रत्येक $E$ ला ऋणेतर मूल्य $\mu(E) \in
\Real$ दिले जाते; हे $X$ चे माप ठरवते. तपशील येथे आवश्यक नाहीत; पण फलन $\mu$ ने गणनीय योगशीलतेचे तत्त्व पाळले पाहिजे: वियुक्त संचांच्या गणनीय युतीकरणाचे माप त्यांच्या स्वतंत्र मापांची बेरीज असते. म्हणजे $X_n$ हे संच वियुक्त असताना $\mu(\bigcup_{n < \omega} X_n) = \sum_{n < \omega}\mu(X_n)$.
संच “अमापनीय” आहे याचा अर्थ तो या निश्चित मापाच्या मापनीय क्षेत्रात येत नाही; कोणतेही माप कधीच देता येत नाही, असा सार्वत्रिक अर्थ नाही. आता आधीचा $\rotationsgroup$ वापरू:

\begin{cor}[विताली]
$\mu$ हे असे माप असू द्या की $\mu(\onesphere) = 1$, आणि $X$ व $Y$ एकरूप असतील, तर $\mu(X) = \mu(Y)$. मग प्रत्येक $\rho \in
\rotationsgroup$ साठी $\funimage{\rho}{C}$ अमापनीय आहे.
\end{cor}

\begin{proof}
विरोधसिद्धीसाठी तसे नाही असे समजा. मग एखाद्या $\sigma \in \rotationsgroup$ आणि $r \in \Real$ साठी $\mu(\funimage{\sigma}{C}) =
r$ असू द्या. कोणत्याही $\rho \in \rotationsgroup$ साठी, $\funimage{\rho}{C}$ आणि $\funimage{\sigma}{C}$ एकरूप आहेत; म्हणून प्रत्येक $\rho \in \rotationsgroup$ साठी $\mu(\funimage{\rho}{C}) = r$.
\olref{vitalicover} आणि \olref{vitalinooverlap} नुसार $\onesphere =
\bigcup_{\rho \in \rotationsgroup}\funimage{\rho}{C}$ हे जोडीजोडीने वियुक्त संचांचे गणनीय युतीकरण आहे. म्हणून गणनीय योगशीलतेनुसार $\mu(\onesphere) = 1$ ही प्रत्येक $\funimage{\rho}{C}$ च्या मापांची बेरीज आहे; म्हणजे
\[
	1 = \mu(\onesphere) = \sum_{\rho \in \rotationsgroup}\mu(\funimage{\rho}{C}) = \sum_{\rho \in \rotationsgroup}r
\]
पण $r = 0$ असल्यास $\sum_{\rho \in \rotationsgroup}r = 0$, आणि $r
> 0$ असल्यास $\sum_{\rho \in \rotationsgroup}r = \infty$.
माप ऋणेतर असल्याने हीच दोन प्रकरणे शक्य आहेत; आणि दोन्हीमध्ये सामान्यीकरणाशी विसंगती मिळते.
\end{proof}

\end{document}
